\chapter{実験結果詳細}
\label{app:details}

\section{再現のための記録}
付録では，本文の流れを妨げずに詳細な情報を残すことができる．
本例では，実験を再現するために記録しておく項目を整理する．
資料の名称だけでなく，実際に使用した版や設定を記載し，本文中の表と結果ファイルの対応が
分かるようにすることが望ましい．ファイルを保存するだけでなく，どの処理で作られたかも説明する．

確認する項目の例を以下に示す．
\begin{enumerate}
  \item 文書集合，検索要求，適合性判定の版を記録する．
  \item 前処理と順位付けに使用した設定を保存する．
  \item 検索要求ごとの出力と評価値を対応付ける．
  \item 除外した要求とその理由を明記する．
  \item 各表を再生成するための集計手順を記載する．
\end{enumerate}

付録の内容は，本文から参照されることで役割が明確になる．
例えば，第\ref{sec:experiments}章の集計方法を再現するために必要な設定を示す，という関係を説明するとよい．
ページ番号を手入力する代わりに参照命令を使用すれば，本文が増減しても参照先を更新できる．
このサンプルでは，付録にアルファベットの番号を付け，節番号にも同じ文字を使用する．

% Demonstration-only break.
\newpage
\section{文書全体の点検}
このページは，付録の通常ページの書式を確認するための例である．
付録の最初のページにはページ番号だけを表示し，続きのページには見出しと横線も表示する．
目次に示される付録のページ番号は，付録の最初のページを指していなければならない．
本文の各章と同様に，付録の前にも奇数ページへ合わせるための空白ページは挿入しない．

図や表を追加すると，本文のページ分割が変化することがある．
これは通常の組版の動作であるが，図題と図が離れていないか，参照先の番号が更新されたか，
大きな表が余白を越えていないかを確認する必要がある．長い見出しについても，
本文とページ上部の見出しの両方で，文字が重なっていないかを点検する．

正式な論文を執筆する際には，例文と架空の数値，架空の文献を削除する．
表示確認のために設けた改ページも取り除き，実際の内容を主ファイルに記述してよい．
表紙の所属，学籍番号，氏名，提出日は実際の情報に置き換える．
文書の体裁が整っていても，内容の正確さが自動的に保証されるわけではないため，両方を確認することが重要である．
